\chapter{Lessons learned}
\label{ch:lessons_learned}

\grokfig{fig_pi_tantum_buonanotte}{$\pi$Tantum a fine giornata,
in pigiama a righe e cuffietta da notte, con candela in mano.}

\epigraph{``Experience is what you get when you didn't get
what you wanted.''}{Randy Pausch, \emph{The Last Lecture}}

\lettrine[lines=3]{D}{opo} un anno e mezzo di sviluppo,
qualche centinaio di esperimenti e una manciata di migliaia
di righe di codice, alcuni insegnamenti stanno emergendo. Li
raccolgo qui per chi voglia costruire un sistema simile, o
estendere $\pi$Tantum, o semplicemente capire dove le
intuizioni iniziali sono state confermate e dove smentite.

\section{Cosa ha funzionato}

\subsection{Modellare HARD/SOFT separatamente}

Gi\`a all'inizio era chiaro che mescolare i due tipi di
vincoli sarebbe stato un disastro. La distinzione netta
``HARD entra come vincolo nel solver, SOFT entra come
penalit\`a nell'obiettivo'' \`e \emph{l'unica} che permette
all'utente di leggere correttamente i risultati. Senza di
essa l'utente non distingue ``il sistema non riesce a
soddisfarlo'' da ``il sistema sta cercando di soddisfarlo
ma trova un compromesso''. La palette 5-stati con codice
colore \`e l'estensione naturale che permette di esprimere
sfumature (preferred, enforced) senza aumentare la
complessit\`a percepita.

\subsection{Decomporre invece di scalare}

L'illusione del primo periodo era che ``CP-SAT scali
abbastanza''. Vero fino a $\sim$50 classi. Sopra, i tempi
esplodono. La decomposizione spettrale (cap.~\ref{ch:metodo_spettrale})
\`e stata l'introduzione che ha cambiato l'orizzonte di
applicabilit\`a: con cluster naturali ben separati, lo
speedup \`e di un ordine di grandezza.

\subsection{DSL parser hand-rolled}

La tentazione iniziale era usare PLY o Lark. Alla fine il
parser hand-rolled \`e risultato pi\`u pulito: 250 righe di
Python, niente dipendenze, controllo esatto degli errori.
Quando aggiungi un built-in nuovo, sai dove toccare. Quando
un utente sbaglia la sintassi, l'errore \`e in italiano e
indica la posizione del token. Le librerie generiche
sarebbero state pi\`u verbose nello stesso codice e meno
``parlanti'' negli errori.

\subsection{Frontend lazy: non fare polling per niente}

Caricare ECharts solo on-demand, fare polling solo se la
tab \`e visibile (\texttt{document.hidden} check), unmount
delle pagine al cambio tab: piccoli accorgimenti che
trasformano un'app ``pesantuccia'' in un'app ``snella'' su
hardware modesto. La scuola che usa $\pi$Tantum ha
tipicamente computer di 5--7 anni; la responsivit\`a percepita
\`e cruciale.

\subsection{Run asincroni con SSE log streaming}

Trasformare ogni operazione lunga (Phase A, Phase B,
Monte Carlo, Hall) in un \emph{run asincrono} con log
streaming via SSE \`e stata una svolta UX. L'utente non aspetta
con la barra di caricamento; legge il log in tempo reale,
capisce cosa sta succedendo, pu\`o killare se serve. La
metafora ``ogni operazione \`e un job che lascia traccia in
\texttt{Run+RunLog+RunTelemetry}'' \`e diventata l'asse
portante dell'architettura.

\section{Cosa non ha funzionato (al primo tentativo)}

\subsection{Pickle come engine I/O}

\sidenote{Pickle \`e \emph{conveniente} ma fragile: cambia la
classe Python e i pickle vecchi non si caricano pi\`u.}

L'idea di scambiare dati fra backend e solver via pickle ha
funzionato per il prototipo ma ha causato dolore quando si
sono raffinate le strutture dati. Ogni cambio di campo nelle
dataclass del solver invalidava i pickle. Soluzione
adottata: \texttt{engine\_io.py} fa da convertitore esplicito
fra DB e dict; il pickle \`e ancora usato come cache, ma con
un campo \texttt{schema\_version} che permette di
riconoscere i pickle vecchi e rigenerarli.

\subsection{Migrazioni hand-rolled vs Alembic}

La scelta consapevole di scrivere migrazioni idempotenti
direttamente in \texttt{db.py} (\emph{senza} Alembic) ha
pro e contro. Pro: zero file di migrazione separati, zero
ambiente da configurare, deploy banale. Contro: con
$\sim 30$ tabelle si comincia a sentire il peso (la funzione
\texttt{\_apply\_lightweight\_migrations()} \`e ormai
lunghissima). Probabile rifattorizzazione in vista per la
versione successiva.

\subsection{Promesse di ottimo globale con metaeuristiche}

\begin{erroricomunibox}
Le metaeuristiche \emph{non} garantiscono l'ottimo. Ho
provato pi\`u volte a ``promettere'' all'utente che la
soluzione finale fosse il meglio possibile. \`E falso, ed \`e
controproducente: l'utente perde fiducia quando vede un
piccolo manual fix migliorare l'obiettivo. La onest\`a
intellettuale paga: ``questa \`e una soluzione di alta
qualit\`a, non garantita ottima'' \`e meglio di ``ottima''.
\end{erroricomunibox}

\subsection{Drag-and-drop senza preview}

Prima versione: drag-and-drop senza preview, conferma
post-hoc se la mossa era HARD-fattibile. Era frustrante.
Versione attuale: preview live con check HARD e delta SOFT
mentre trascini. Adesso l'utente impara giocando: vede
\emph{prima} che la mossa \`e infattibile e capisce
\emph{perch\'e}. La differenza in usabilit\`a \`e enorme.

\subsection{Sottostimare i tempi su \texttt{superhuge}}

Il profilo \texttt{superhuge} (80 classi, 16 sezioni $\times$
5 anni) era pensato come ``stress test estremo'', non come
caso d'uso reale. Si \`e scoperto che esistono \emph{davvero}
istituti italiani con queste taglie (poli scolastici, IIS).
Phase A su \texttt{superhuge} consuma 5 minuti --- quasi
accettabile, ma pi\`u del previsto. Le ottimizzazioni di
Phase A (decomposizione di matching, MIP-based) sono in
roadmap.

\section{Quello che ancora si pu\`o migliorare}

\subsection{DSL editor avanzato}

\sidenote{Tutta la roadmap di DSL editor (squiggle, fuzzy,
quick-fix, hover docs, NL preview) \`e descritta nel
cap.~\ref{ch:metodo_dsl}.}
Ad oggi il campo DSL \`e una semplice \texttt{textarea}
senza assistenza. La specifica completa di un editor con
live validation, suggerimenti fuzzy, quick-fix e hover docs
esiste; manca il tempo di implementarla. \`E la priorit\`a
numero uno per il prossimo ciclo.

\subsection{Documentazione viewer in-app}

Tutta questa documentazione esiste come PDF e markdown nel
repository, ma non \`e accessibile dall'interno dell'app.
Idealmente: una tab \texttt{/docs} con albero a sinistra
(per categoria) e contenuto a destra (markdown reso live),
con search globale.

\subsection{Wizard per casi d'uso comuni}

Casi ricorrenti come ``crea un corso di recupero'' o
``aggiungi una compresenza'' richiedono al coordinatore di
sapere quali campi compilare e in quale tab. Sarebbero un
buon target per dei \emph{wizard} guidati che fanno tutto
in pochi click.

\subsection{Multi-tenant}

Ad oggi $\pi$Tantum \`e mono-utente / mono-scuola: una
istanza per scuola. Per un servizio cloud serve gestione
utenti, isolamento dei dati, billing. \`E un'estensione
significativa che richiede di ripensare il modello dati e
l'autenticazione.

\subsection{Mobile responsivo}

L'interfaccia \`e desktop-first. Per gestire le supplenze
quotidiane dal telefono servirebbe un layout responsive. Non
\`e solo un fatto di CSS: la matrice 6$\times$6 della pagina
``Assenze'' non ci sta su uno schermo da 5 pollici.

\section{Una riflessione finale}

Costruire $\pi$Tantum mi ha confermato un'intuizione che
all'inizio era solo una sensazione: \emph{i sistemi
complessi sopravvivono se sono onesti con l'utente}. Onesti
nel dire ``questo non posso farlo'', onesti nel mostrare
\emph{quanto} una scelta costa, onesti nel non promettere
ottimi che non sono garantiti. L'utente perdona la lentezza
se sa perch\'e. L'utente apprezza il pre-check di Hall che
dice ``ferma, mancano 3 ore di matematica, prima sistemale''
infinitamente di pi\`u del solver che dopo 10 minuti
restituisce ``no soluzione''.

In fondo \`e la stessa lezione che si impara ogni volta che
si lavora con altri umani: la fiducia si costruisce un
piccolo segnale alla volta. Anche un software \`e fatto di
piccoli segnali.

\pullquote{I sistemi complessi sopravvivono se sono onesti con l'utente.}{Lessons learned, $\pi$Tantum 2026}
